% Pursuing Stacks (cotes 134-2 to 134-8): archive cote 134-4, site page 40 (PDF page 41), typescript page 128.
% typescript. Draft compared: ps3.tex:1998-2061.
% AI-transcribed by Claude Opus 5.5 (claude-opus-5-5, adaptive thinking, effort high), Message Batches API, batch ps2 (msgbatch_01NVwr1Yuo8qJf5rykeGBXPK), 2026-11-12; not yet checked by eye.
\scanpage{134-4/40}
\mspage{128}

structure, and the two are associated. This comes from the fact that
both compositions of the functors (2),(3) must be the identity functors
(being compatible with the ``under M'' structure), hence \hand{$\alpha$} and \hand{$\beta$} are isomor%
phisms, which shows both R \hand{$\in$} Q$_{\rm o}$ and W \hand{$\in$} L$_{\rm o}$ , and that R and L are associa-%
ted. In terms of the set theoretic situation (1), this may be described
by using the order relations on Q and L, and the fact that r and s are
monotone maps, and satisfy moreover
(*) \qquad $sr(q)\hand{$\leq$} q$ \;,\; $rs(\hand{$\ell$})\hand{$\leq$}\hand{$\ell$}$ \quad (any $q\hand{$\in$}Q$, $\hand{$\ell$}\hand{$\in$}L$) \;,
which implies that the set C$_{\rm o}$ of pairs $(q\unsure{_{,}}\,, \hand{$\ell$})$ of associated elements
of Q$_{\rm o}$,L$_{\rm o}$ can be described also as
\[ C_{o} = \hand{$\{$}(q,\hand{$\ell$})\hand{$\in$} Q\text{x}L \hand{$\mid$} \hand{$\ell$}\hand{$\leq$}r(q),\quad q\hand{$\leq$} s(\hand{$\ell$})) \qquad . \]

Thus it does'nt seem evident to get a homotopy structure on M, just
starting with an\xout{y} R \hand{$\in$} Q or \xout{\unsure{any}} a W \hand{$\in$} L, without assuming beforehand that
R is a homotopy relation, or W a homotopism structure. The condition
$C_{1,2}$ on page 115 may be viewed as a condition on a pair (R,W)\hand{$\in$} QxL,
and it clearly implies
\[ R \hand{$\leq$} s(W) \quad ; \]
if we assume moreover \quad W=r(R) \quad we get \quad R\hand{$\leq$} s(r(R)) \quad and hence, in view
of the first inequality (*) above ,
\[ R = sr(R) \quad \text{i.e.} \quad R \hand{$\in$} Q_{o} \quad , \]
i.e. R is a homotopy relation, as asserted somewhat quickly yesterday.

The only standart way for getting homotopy structures in a general
category M which I \xout{xxx} can see by now , is in terms of an arbitrary
set \hand{$\Sigma$}$_{h}$ of intervals in M (assuming only that M admits finite products).
As soon as the focus gets upon intervals for describing homotopy structu-%
res, the situation becomes typically non-autodual - in contrast to
Quillen's autodual treatment of homotopy relations (via``left'' and ``right''
homotopies, involving respectively ``cylinder'' and ``path'' - objects). \xout{Thi}
This is in keeping with the \inserted{(highly non auto-dual} axiom\deleted{\unsure{s}} on universal disjoint \uline{sums in M,} which
we finally introduced by the end of yesterday's \uline{reflection.}\hand{$|$} To come
back however upon the relationships between the four basic ``homotopy
notions'' introduced in yesterday's notes, I would now rather symbolize
these relations in the following diagram of maps between sets

\deleted{\illegible{1}}

\margin{40}
